\section{Removing the longitudinal seed controls}
\label{sec:seed-simulation}

Theorem~\ref{thm:classical-synthesis} produces a finite word that uses a Gaussian
longitudinal field in addition to the cost and mixer.  We now remove
that extra control while preserving its limiting tree energy.  This is
the last compilation step: the fixed-word transfer then applies to the
resulting ordinary QAOA word.

A \emph{seeded word} is a finite word in the cost, the mixer, and
\(\exp(-\ii\sigma\sum_i Z_i s_i)\), where the \(s_i\) are independent
standard real Gaussians.  The same Gaussian family is used at every
seed gate.  The Gaussian variables are averaged in evaluating the word;
equivalently, they are multiplication operators on independent pointed
Gaussian spaces.  All the angles, including \(\sigma\), are deterministic.
We fix this entire seeded word before choosing any of its approximants.
Its limiting tree energy is \(\cE(v_Z)\), where \(v_Z\) is its
root insertion vector. The seeded root update in
Lemma~\ref{lem:ordinary-updates} gives this vector and the uniform
occupation bound, so Proposition~\ref{prop:energy} identifies its
edge-energy limit with \(\cE(v_Z)\). This uses the seeded tree
calculus, not an extension of the ordinary-word SK transfer.

\begin{theorem}[Removal of seed controls]
\label{thm:seed-simulation}
For every finite seeded word and every \(\epsilon>0\), there is a finite
ordinary cost--mixer word whose limiting tree energy differs from that of
the seeded word by less than \(\epsilon\).  Its angles are independent of the
degree, the system size, and the disorder.  In particular, the approximation
uses the thermodynamic limit at a fixed word.
\end{theorem}

The proof uses two parameters, chosen in order.  We first choose a small
fixed \(b>0\).  With \(b\) fixed, \(r\downarrow0\) creates a component of
amplitude \(r\) with a Weyl displacement of order \(r^{-1}\).  Short cost pulses can
selectively rotate this displaced component while their first-order action
cancels on the component of order-one amplitude.  A cost echo of duration
\(r^{-1}\) converts the change of the small component into a finite
longitudinal field.  All approximation statements below are statements in
the already established fixed-word tree limit.

We first establish estimates for displaced Fock vectors and prepare
an exact two-component state.  Finite selective spin rotations then let
us propagate its small displaced component through the requested word.
The echo calculation produces the seed field; a separate root-shift
calculation identifies its energy.  Finally, a finite generic
perturbation supplies all the nonzero projections required to choose
the selective rotations.

\subsection{Uniform estimates for packets}

The large displacements will not have bounded occupation moments.
We therefore keep each displacement separate from a tail vector
whose moments are controlled.  Write \(N\) for the number operator in
the \emph{top} child Fock space and put
\[
 \mathcal S=\bigcap_{k\geq0}\operatorname{Dom}(1+N)^k,
 \qquad \norm{\eta}_k=\norm{(1+N)^k\eta}.
\]
The spin factor is included in this notation.  No weight is imposed inside
a child-mode direction.  A bounded family of inputs means a family
\(\eta_r\in\mathcal S\) such that
\(\sup_r\norm{\eta_r}_k<\infty\) for every \(k\).  Uniform errors below are
uniform on every such family.  More explicitly, every estimate used below
depends on finitely many of these seminorms; its constants may depend on
the fixed finite pulse list and on \(b\), but not on \(r\) or on the input.

Let
\[
 W(d)=e^{\ii q(d)},\qquad
 W_P(d)=e^{\ii Pq(d)},\qquad P^2=\1,
\]
where \(P\) is a root Pauli matrix.  The Weyl identities are
\begin{equation}
\label{eq:seed-weyl}
 W(d)W(e)=e^{-\ii\im\braket d e}W(d+e),
 \qquad
 W(d/r)^*q(v)W(d/r)
 =q(v)-\frac{2}{r}\im\braket v d.
\end{equation}
A packet with scaled center \(d_r\) is \(W(d_r/r)\eta_r\), with
\(\eta_r\) a bounded family.  A finite sum of packets is allowed, as are
bounded root spin matrices in their tails.  Scalar phases of modulus one
may depend on \(r\) and need not converge.

\begin{samepage}
\begin{lemma}[Packet estimates]
\label{lem:seed-packet-estimates}
The following estimates are uniform for bounded families of tails.
\begin{enumerate}
\item Displacements with bounded directions preserve all the seminorms
\(\norm\cdot_k\), with bounds depending only on the direction norm and
\(k\).  For \(\norm d\leq C\),
\[
 W(rd)\eta=\eta+\ii r q(d)\eta+O(r^2)\norm\eta_1
 \quad\hbox{in norm}.
\]
If \(d_r-e_r\to0\) in norm and both directions are bounded, then
\((W(d_r)-W(e_r))\eta_r\to0\) in norm.  These assertions also hold
for \(W_P\).
\item If \(\norm{d_r-e_r}\geq\kappa>0\), then for every integer \(M\),
\begin{equation}
\label{eq:seed-packet-separation}
 \abs{\braket{W(d_r/r)\eta_r}{W(e_r/r)\xi_r}}=O(r^M).
\end{equation}
The conclusion is unchanged by inserting a bounded root spin matrix.
\item For bounded \(d_r,v_r\) and fixed \(\tau\), conjugating the short
probe \(e^{-\ii r\tau Pq(v_r)}\) through a packet gives
\begin{equation}
\label{eq:seed-packet-probe}
 W(d_r/r)^*e^{-\ii r\tau Pq(v_r)}W(d_r/r)
 =e^{-\ii\tau a_rP}e^{-\ii r\tau Pq(v_r)},
 \qquad a_r=-2\im\braket{v_r}{d_r}.
\end{equation}
Consequently a fixed finite list of short probes preserves its scaled
centers and bounded centered-tail moments, without any smallness
assumption on \(a_r\).
\end{enumerate}
\end{lemma}
\end{samepage}

\begin{proof}
The first assertion follows from the creation--annihilation estimates in
Lemma~\ref{lem:weyl-bounds}, or by differentiating a Weyl displacement on
\(\mathcal S\).  In the Taylor remainder one uses
\(\norm{q(d)^2\eta}\leq C\norm d^2\norm\eta_1\).
The Weyl product identity then gives the assertion for two nearby
directions.  Resolving \(P\) into its two spectral projections proves the
spin-controlled statements.

For the second assertion, the Weyl product reduces the question to a
matrix coefficient of \(W(f_r/r)\), where \(\norm{f_r}\geq\kappa\).
Choose the unit quadrature
\(Q_r=q(\ii f_r/\norm{f_r})\).  Equation~\eqref{eq:seed-weyl}
translates \(Q_r\) by \(2\norm{f_r}/r\).  Split each vector using the
spectral projection of \(Q_r\) onto
\([-\norm{f_r}/(2r),\norm{f_r}/(2r)]\).  The two translated central
intervals are disjoint.  The complementary pieces have norm bounded by
\(C_Mr^M\norm\eta_{k_M}\), by the spectral theorem and the uniform
estimate \(\norm{Q_r^M\eta}\leq C_M\norm\eta_{k_M}\).
This proves~\eqref{eq:seed-packet-separation}.  Root spin matrices commute
with these projections.  Finally,~\eqref{eq:seed-packet-probe} is exactly
the second Weyl identity; its right-hand side is a spin rotation followed
by a displacement of size \(O(r)\).  Both preserve the required tail
bounds.  Notice that this last argument precedes, and does not use, any
claim that the probe acts almost trivially on an ordinary packet.
\end{proof}

We use exact local unitaries and exact insertion directions throughout.
At a given stage, if the local unitary is \(u_r\), put
\[
 v_{P,r}=u_r^*Pu_r\vac.
\]
All the \(v_{P,r}\) have norm one.  A toggled cost probe about an axis
\(P\) in the \(ZY\)-plane is
\(e^{-\ii r\tau Pq(v_{P,r})}\).  It is an ordinary word: conjugate a
cost pulse by two opposite mixer pulses.  The direction is the insertion
at the start of that probe; it is constant during the probe itself.

\subsection{Initialization and the exact two-component decomposition}

We next prepare the two components that the selective probes will distinguish.
Their decomposition is exact, even though they need not be orthogonal.
This exactness will keep the special component's amplitude equal to
\(r\) through all later unitary updates.

Fix once and for all
\(\beta_\star=\pi/4\), and put \(R_\star=e^{-\ii\beta_\star X}\).
This choice makes the two eventual external Gaussian directions
orthogonal.  For \(0<r<1\), set
\[
 \delta_r=\arcsin r,\qquad \gamma_r=b/r,\qquad e=Z\vac,
 \qquad
 \ell_r=Y e^{-2\ii\gamma_r Zq(e)}\vac.
\]
Apply, in order, cost \(\gamma_r\), mixer \(\delta_r\), cost
\(-\gamma_r\), and mixer \(\beta_\star\).  Immediately before the
negative cost the insertion is
\[
 v_{2,r}=\cos(2\delta_r)e+\sin(2\delta_r)\ell_r.
\]
Define \(d_{0,r}=\gamma_r(v_{2,r}-e)\).  The resulting local unitary
has the \emph{exact} decomposition
\begin{equation}
\label{eq:seed-initial-split}
 u_r=a_{0,r}U_{O,r}+a_{1,r}U_{S,r},\qquad
 U_{O,r}=R_\star W_Z(d_{0,r}),\qquad
 U_{S,r}=R_\star XW_Z(-2\gamma_re-d_{0,r}),
\end{equation}
where
\[
 a_{0,r}=\cos\delta_r\,e^{\ii\vartheta_r},\quad
 a_{1,r}=-\ii r e^{-\ii\vartheta_r},\quad
 \vartheta_r=\gamma_r^2\im\braket{v_{2,r}}e.
\]
Indeed, this is obtained by expanding \(e^{-\ii\delta_rX}\) and applying
the Weyl product identity.  Since
\[
 \braket e{\ell_r}=-\ii e^{-2\gamma_r^2},\qquad
 d_{0,r}=-2br\,e+2b\sqrt{1-r^2}\,\ell_r,
\]
we have \(\vartheta_r\to0\) faster than every power of \(r\), and
\(\norm{d_{0,r}-2b\ell_r}\to0\).

At every later stage retain the exact decomposition obtained by applying
all subsequent local updates to the two operators in
\eqref{eq:seed-initial-split}.  Denote its actions on an input \(\eta\)
by \(\Psi_{O,r}(\eta)\) and \(\Psi_{S,r}(\eta)\).  In particular,
\begin{equation}
\label{eq:seed-exact-special-norm}
 u_r\eta=\Psi_{O,r}(\eta)+\Psi_{S,r}(\eta),\qquad
 \norm{\Psi_{S,r}(\eta)}=r\norm\eta.
\end{equation}
No orthogonality is asserted or needed in this exact split.
The echo will depend on the change in the root \(Z\)-insertion when
the special component is flipped.  We therefore record the insertion
carried by that component, divided by its exact amplitude \(r\).
Let \(\psi_{S,r}=\Psi_{S,r}(\vac)\), and define at initialization
\begin{equation}
\label{eq:seed-address-direction}
 w_r=u_r^*Z\psi_{S,r}/r.
\end{equation}
Thus \(\norm{w_r}=1\) exactly.  The next lemma separates this
direction and the initialization direction \(\ell_r\) from all
bounded families of top-number-smooth vectors.

\begin{lemma}[The initial high directions]
\label{lem:seed-high-directions}
The vectors \(\ell_r,w_r\) are centered, have norm one, and satisfy
\[
 \braket{\ell_r}{w_r}\longrightarrow0,
 \qquad
 \norm{\1_{N\leq M}\ell_r}+
 \norm{\1_{N\leq M}w_r}\longrightarrow0
 \quad\hbox{for every fixed }M.
\]
They consequently have vanishing inner products with every bounded family
of top-number-smooth vectors.
\end{lemma}

\begin{proof}
Centering follows from the global spin-flip symmetry; every initialization
gate is even and the inserted \(Y\) or \(Z\) is odd.  The norm statements
follow from unitarity.  Because
\(R_\star^*ZR_\star=Y\) and \(YX=-\ii Z\), expansion of
\eqref{eq:seed-address-direction} gives
\[
 w_r=-Z W_Z(-2\gamma_re-2d_{0,r})\vac+o(1),
\]
where the omitted special--special term has norm \(r\), and the remaining
scalar phase tends to one.  The directions \(d_{0,r}\) are bounded.
Resolving the root spin shows that both displayed high vectors are sums
of coherent packets whose displacements have norm tending to infinity.
Their projections onto a fixed number sector therefore vanish, directly
from the coherent-state expansion.  At a fixed output \(Z\)-spin the
centers in \(\ell_r\) and in the displayed expression for \(w_r\) differ
by \(4\gamma_re+O(1)\); their overlap tends to zero by
Lemma~\ref{lem:seed-packet-estimates} (or the vacuum Weyl formula).
Finally, split a bounded family at \(N\leq M\).  Its complementary norm
is uniformly \(O(M^{-k})\), while its fixed-sector pairing tends to zero.
First take \(r\downarrow0\) and then \(M\to\infty\).
\end{proof}

The special part of~\eqref{eq:seed-initial-split}, on \emph{any} bounded
input family, has four packets.  Index them by the spin \(z=\pm1\)
before \(X\) and by the spin \(y=\pm1\) after \(R_\star\).  Their
scaled centers are \(-2bz e\), and their output spins are \(y\).
The bounded displacement \(-z d_{0,r}\) is part of the tail.  In
particular, the center list is independent of the input; only the tails
depend on it.

\subsection{A finite ensemble of spin rotations}

On the vacuum, the exact split has an ordinary component with bounded top-number
moments and a special component of norm \(r\) carried by four displaced
packets.  A short cost pulse acts on each displaced vector mainly as a
spin rotation.  We need a finite list of such pulses that gives the same
rotation at every special center, while its first-order action on an
undisplaced vector cancels.  The two zero sums in the next lemma impose
that cancellation separately for the two available spin axes.

The selective-rotation argument uses the Lie-bracket and polynomial-
approximation approach to ensemble control of Li and
Khaneja~\cite[Section~II, Examples~1--2]{LiKhaneja2005}. The finite-family
and exact zero-clock requirements needed here are proved below.

\begin{lemma}[Selective rotations with zero clocks]
\label{lem:seed-ensemble-control}
Let \(h_1,\ldots,h_m\) be a finite list of nonzero vectors in \(\R^2\).
There are two perpendicular axes \(P_1,P_2\) in the \(ZY\)-plane with
the following property.  Set \(c_{j,a}=h_j\cdot P_a\), using the
coordinates \((Z,Y)\).  Given any \(\theta\in\R\), there is a finite
list \((a_k,\tau_k)\), with \(a_k\in\{1,2\}\), such that
\begin{equation}
\label{eq:seed-zero-clocks}
 \sum_{k:a_k=a}\tau_k=0\quad(a=1,2),
 \qquad
 \prod_k e^{-\ii\tau_k c_{j,a_k}P_{a_k}}
 =e^{-\ii\theta X}\quad(1\leq j\leq m).
\end{equation}
The product is in chronological order, with later factors on the left.
The assertion is exact, not merely an assertion about closure.
\end{lemma}

\begin{proof}
Choose the axes so that every \(c_{j,a}\neq0\), and so that the pairs
\((c_{j,1}^2,c_{j,2}^2)\) are distinct unless the original vectors are
equal or opposite.  Such axes exist: each prohibited equality or zero is
a nontrivial trigonometric equation in the axis angle, apart from the
stated equal-or-opposite exception.  Delete duplicate square pairs; on
each deleted pair the product \(c_{j,1}c_{j,2}\) is the same.

In \(G=SU(2)^m\times\R^2\), take the two generators having spin
components \(-\ii c_{j,a}P_a\) and clock component the \(a\)-th
coordinate vector.  Their brackets have zero clocks.  A first bracket
has \(\ii X\)-coefficient equal, up to a common nonzero real constant,
to \(c_{j,1}c_{j,2}\).  Repeated double brackets with the first or
second generator multiply this coefficient by \(c_{j,1}^2\) or
\(c_{j,2}^2\), again up to nonzero common constants.  Thus the generated
Lie algebra contains the clock-zero vector with coefficients
\[
 c_{j,1}c_{j,2}F(c_{j,1}^2,c_{j,2}^2)\,\ii X
\]
for every real polynomial \(F\).  Polynomial interpolation on the finite
set of square pairs makes all these coefficients equal to one.

For completeness, exponentials of this Lie-algebra vector are reachable
by finite products of the two original one-parameter groups.  Let
\(\mathcal R\) denote the subgroup of such finite products, and let
\(\mathfrak g_0\) be the Lie algebra generated by the two generators.
The span of their conjugates by elements of \(\mathcal R\) is invariant
under each of the two adjoint one-parameter groups.  Differentiation
shows that it contains all iterated brackets and hence equals
\(\mathfrak g_0\).  Choose finitely many of these conjugates forming a
basis.  The product map of their one-parameter groups has surjective
differential at zero, so its image contains a neighborhood of the
identity in the connected analytic subgroup with Lie algebra
\(\mathfrak g_0\).  Each factor is itself a reachable conjugate.
Subdivide any one-parameter path into finitely many pieces in this
neighborhood.  The whole exponential is therefore in \(\mathcal R\).
Applying this to the common \(X\)-generator with zero clocks gives
\eqref{eq:seed-zero-clocks}.
\end{proof}

The preceding lemma treats fixed real coefficients.  In the circuit,
the coefficients come from actual insertion vectors and change during
the probes.  The next estimate controls that change without replacing
the insertion vectors by approximants before a singular rescaling.

\begin{lemma}[Accuracy of selective probes]
\label{lem:seed-selective-accuracy}
Consider a stage whose vacuum state \(u_r\vac\) has a finite packet
expansion, with norm-\(o(r)\) remainders and uniformly bounded scaled
centers.  Suppose the ordinary component has order-one amplitude and
the special component has amplitude \(O(r)\).  All designated tails have
bounded top-number moments.  Suppose also that one of the following
ordinary geometries holds.
\begin{enumerate}
\item At a standard stage all ordinary centers are zero.
\item At a huge-cost stage the ordinary part is
\(\psi_++\psi_-\), with \(Z\psi_y=y\psi_y\) and centers
\(-tyv_r\), where \(t\neq0\) and \(v_r=u_r^*Zu_r\vac\).
Every special center \(d_{j,r}\) obeys
\(\lim\im\braket{v_r}{d_{j,r}}\neq0\).
\end{enumerate}
Assume the limiting special projection pairs
\[
 h_j=\lim_{r\downarrow0}-2
 \bigl(\im\braket{v_{Z,r}}{d_{j,r}},
       \im\braket{v_{Y,r}}{d_{j,r}}\bigr)
\]
exist and are nonzero.  Use the axes and finite pulse list of
Lemma~\ref{lem:seed-ensemble-control} for these pairs, replacing its
spin generators by actual toggled probes of angles \(r\tau_k\).
Then the net probe action is the identity up to \(o(r)\) on every
ordinary packet, and is the prescribed \(X\)-rotation up to \(o(1)\)
on every normalized special packet.  The ordinary assertion holds also
after any root spin matrix has been inserted in its tail.  These
statements hold uniformly for bounded families of tails sharing the
stated center list, with the probes chosen solely from the vacuum data.
\end{lemma}

\begin{proof}
Let \(V_j\) be a prefix of the short probes, and write
\(v_{P,r}^{(j)}=(V_ju_r)^*P(V_ju_r)\vac\).  Center preservation and
all centered-tail bounds follow first from
Lemma~\ref{lem:seed-packet-estimates}, without estimating the spin
phases.  At a standard stage every ordinary probe has zero center phase.

At a huge-cost stage there is the exact identity
\begin{equation}
\label{eq:seed-hermitian-diagonal}
 \braket{v_{P,r}^{(j)}}{v_r}
 =\braket{PV_j\psi}{V_jZ\psi},\qquad \psi=u_r\vac.
\end{equation}
Group the ordinary packets by their original sign \(y\).  The diagonal
term of sign \(y\) on the right-hand side is
\(y\braket{PV_j\psi_y}{V_j\psi_y}\), which is real.  This remains true
although the probes may rotate the spin within that group.  The two
ordinary centers are separated by \(2|t|\), since \(\norm{v_r}=1\).
An ordinary center and a special center are also separated: their
imaginary projections against \(v_r\) differ by a quantity bounded away
from zero.  Lemma~\ref{lem:seed-packet-estimates} makes all these cross
terms smaller than any fixed power of \(r\).  Terms containing only
special packets are \(O(r^2)\), and the remainder terms are \(o(r)\).
Consequently, at \emph{every} intermediate probe prefix,
\begin{equation}
\label{eq:seed-huge-imaginary}
 \im\braket{v_{P,r}^{(j)}}{v_r}=o(r).
\end{equation}
More quantitatively, let \(\varepsilon_r=o(r)\) bound the norm of the
discarded vacuum remainder at the start of the probe sequence.  For
every fixed \(M\), the preceding argument gives
\begin{equation}
\label{eq:seed-huge-imaginary-modulus}
 \max_j\abs{\im\braket{v_{P,r}^{(j)}}{v_r}}
 \leq C_M\bigl(r^2+\varepsilon_r+r^M\bigr).
\end{equation}
The constant depends on the finite pulse list, a positive lower bound
on center separations, and finitely many designated vacuum-tail
seminorms.  The exact probe prefixes preserve the norm of the discarded
remainder; no number-operator estimate is imposed on it.
This derivation used only center preservation, so there is no circular
small-phase assumption.

By~\eqref{eq:seed-packet-probe}, an ordinary center \(-tyv_r\) therefore
has spin phase \(o(r)\) in every probe.  In either geometry a probe acts
as identity plus \(O(r)\) on the ordinary part and on that part with any
spin matrix inserted.  On the special part the norm is \(O(r)\), so
unitarity suffices.  Using
\[
 v_{P,r}^{(j+1)}-v_{P,r}^{(j)}
 =(V_ju_r)^*\bigl(V^*PV-P\bigr)(V_ju_r)\vac
\]
for the next probe \(V\), or equivalently estimating
\((PV-VP)(V_ju_r)\vac\), gives
\begin{equation}
\label{eq:seed-probe-direction-change}
 \norm{v_{P,r}^{(j)}-v_{P,r}^{(0)}}=O(r)
\end{equation}
for the fixed finite list.

After removing an ordinary center, expand the displacement in each
probe to first order.  Equation~\eqref{eq:seed-huge-imaginary} supplies
only an \(o(r)\) additional phase in the huge-cost geometry.  Replacing
\(v_{P,r}^{(j)}\) by its starting value in the first-order displacement
costs \(O(r^2)\), by~\eqref{eq:seed-probe-direction-change}.
The first-order sum is therefore
\[
 -\ii r\sum_{a=1}^2
 \left(\sum_{k:a_k=a}\tau_k\right)P_aq(v_{P_a,r}^{(0)})=0.
\]
Here is a uniform remainder bound for this cancellation.  For a fixed
ordinary center \(d_r\), let \(\mathcal B_r(d_r)\) be the net probe
action conjugated by \(W(d_r/r)\), and put
\[
 \Delta_r=\max_j\norm{v_{P_j,r}^{(j)}-v_{P_j,r}^{(0)}},\qquad
 \alpha_r(d_r)=\max_j
       2\abs{\tau_j\im\braket{v_{P_j,r}^{(j)}}{d_r}}.
\]
Then for every designated centered tail \(\eta\),
\begin{equation}
\label{eq:seed-uniform-probe-bound}
 \norm{(\mathcal B_r(d_r)-\1)\eta}
 \leq C_K\left[
       (r^2+r\Delta_r)\norm\eta_1
           +\alpha_r(d_r)\norm\eta\right].
\end{equation}
Indeed, remove each spin phase by its operator-norm bound and replace
each displacement direction by its starting value.  The latter
replacement costs at most \(C r\Delta_r\norm\eta_1\) through the
finite product.  In the product with starting directions, hold those
directions fixed at their given \(r\)-values and introduce an auxiliary
pulse strength \(\lambda\).  Its first derivative at \(\lambda=0\)
is zero, by the clocks.  Its second derivative is a
finite sum controlled by
\(\norm{q(v)q(w)\eta}\leq C\norm v\norm w\norm\eta_1\);
the intervening displacements preserve this seminorm uniformly.
Evaluating the Taylor remainder at \(\lambda=r\) gives
\(C_Kr^2\norm\eta_1\); no differentiability of the directions in
\(r\) is assumed.
The constant \(C_K\) depends only on the fixed pulse list; all actual
insertion directions have norm one.  At a standard center
\(\alpha_r=0\).  At a huge-cost center~\eqref{eq:seed-huge-imaginary-modulus}
gives \(\alpha_r\leq C(r^2+\varepsilon_r+r^M)=o(r)\).
Together with \(\Delta_r=O(r)\), this proves the claimed uniform
\(o(r)\) estimate using the single tail seminorm \(\norm\eta_1\).
It applies unchanged to a root-spin-inserted tail.  If the input also
contains a discarded norm remainder \(\rho_r\), its contribution to
the net identity error is at most \(2\norm{\rho_r}\), by unitarity.
No field operator is applied to such a remainder in this argument.

On a special packet the change of its phase coefficient during the
sequence is \(O(r)\), by~\eqref{eq:seed-probe-direction-change} and the
boundedness of its scaled center.  The small displacements contribute
\(O(r)\) on its tail, and its starting phase coefficients converge to
the prescribed \(c_{j,a}\).  The resulting finite spin product thus
converges to~\eqref{eq:seed-zero-clocks}.  Multiplication by the packet's
amplitude \(O(r)\) makes its error \(o(r)\).
All quantities defining the probes came from the vacuum.  The estimates
after removing a center, however, use only the tail seminorms.  They
therefore apply to every bounded input family with the same center list.
\end{proof}

\subsection{The packet invariant and the singular echo}

We now apply the selective estimates to the exact initialized
decomposition and propagate it through costs, mixers, and simulated
seed gates.  The required statements include
moving input vectors: later insertion estimates must be allowed to use
the approximating circuit itself as an input.

We call the times between simulated gates \emph{standard stages}.  A
second kind of stage occurs immediately after the forward cost in a
seed echo.  The induction will maintain the following assertions, with
the same center list for every bounded input family.
\begin{enumerate}
\item At a standard stage there is a unitary \(A_r\), a finite product
of root spin rotations and spin-controlled Weyl displacements with
uniformly bounded directions, and a scalar \(\zeta_r\) of modulus one,
such that
\begin{equation}
\label{eq:seed-low-ordinary}
 \Psi_{O,r}(\eta_r)=\zeta_r A_r\eta_r+o(r).
\end{equation}
The operators \(A_r,A_r^*\) preserve all top-number seminorms uniformly.
\item The special part has a finite packet expansion up to \(o(r)\),
with amplitude \(r\), bounded tails, scaled centers
\begin{equation}
\label{eq:seed-special-centers}
 d_{zy,r}=-2bz e+yD_r,\qquad z,y\in\{-1,1\},
\end{equation}
and output spin \(y\).  Coincident centers cause no problem.  The vector
\(D_r\) is a finite sum of actual earlier insertion directions with
fixed bounded coefficients.  Each such direction is within \(O(r)\)
of a bounded family in \(\mathcal S\).
\item At a standard stage the special insertion is transported:
\begin{equation}
\label{eq:seed-transport}
 u_r^*Z\psi_{S,r}/r=w_r+o(1),
\end{equation}
where \(w_r\) is the initial vector
\eqref{eq:seed-address-direction}.  Between a special flip and its
undoing the right-hand side instead has a minus sign.
\end{enumerate}
Here and below the norm remainders need not have bounded number moments.
They are kept as separate remainder vectors and carried through later
unitaries by their norm bounds.  Only the designated tails require
moment estimates.

The initialization satisfies this invariant with \(A_r=U_{O,r}\),
\(D_r=0\), and \(\zeta_r=e^{\ii\vartheta_r}\): its ordinary amplitude
differs from one by \(O(r^2)\).  Equation~\eqref{eq:seed-transport} is
initially exact.  The following observation will be used repeatedly.

\begin{lemma}[Insertion approximation from forward action]
\label{lem:seed-low-insertions}
If~\eqref{eq:seed-exact-special-norm} and
\eqref{eq:seed-low-ordinary} hold uniformly on bounded input families,
then
\begin{equation}
\label{eq:seed-insertion-O-r}
 \norm{u_r^*Pu_r\vac-A_r^*PA_r\vac}=O(r)
 \qquad(P=Y,Z).
\end{equation}
In particular the actual insertion directions have bounded smooth
approximants.  No approximation to \(u_r^*\) on a high packet is needed.
\end{lemma}

\begin{proof}
Put \(\widehat v_{P,r}=A_r^*PA_r\vac\).  This is a bounded smooth
family.  Unitarity and \(A_r\widehat v_{P,r}=PA_r\vac\) give
\[
 \norm{u_r^*Pu_r\vac-\widehat v_{P,r}}
 =\norm{Pu_r\vac-u_r\widehat v_{P,r}}
 \leq\norm{P(u_r-\zeta_rA_r)\vac}
    +\norm{(u_r-\zeta_rA_r)\widehat v_{P,r}}.
\]
Each term is \(O(r)\) by the exact special norm and the ordinary
approximation.  The possibly nonconvergent scalar phase cancels.
\end{proof}

For this induction, assume that the nonzero projection hypotheses of
Lemma~\ref{lem:seed-selective-accuracy} hold at every required probe
start.  We call these the \emph{visibility conditions}.  The final
subsection proves that one perturbation of the finite target word
ensures all of them and then constructs the probes in order.

An unscaled cost uses the actual insertion \(v_{Z,r}\) and simply
left-multiplies \(A_r\) by \(e^{-\ii cZq(v_{Z,r})}\).
On a special packet with spin \(y\), it adds a bounded displacement to
the tail and a scalar phase, leaving its scaled center and spin
unchanged.  It commutes with \(Z\), so it preserves the transport
identity exactly.

To implement a bulk mixer \(R_\theta=e^{-\ii\theta X}\), first use
selective probes to apply \(R_\theta^{-1}\) on the special packets,
then apply the global mixer \(R_\theta\).  On ordinary packets the
first step is identity to \(o(r)\), so the low update is \(R_\theta\).
On special packets the two spin actions cancel to \(o(1)\).  The same
is true on their \(Z\)-inserted tails.  Thus the special center and
spin return to their initial values, and transport is preserved to
\(o(1)\) after division by the special amplitude \(r\).

For a longitudinal seed gate use a fixed \(t\neq0\) and the four
operations
\begin{equation}
\label{eq:seed-echo-list}
 \text{cost }t/r,\qquad
 \text{selective }(-\ii X),\qquad
 \text{cost }-t/r,\qquad
 \text{selective }(\ii X).
\end{equation}
The next lemma verifies both the singular cancellation and its error
scale.

\begin{lemma}[One echo]
\label{lem:seed-one-echo}
Under the invariant and the visibility conditions, the echo
\eqref{eq:seed-echo-list} preserves the invariant.  Up to a scalar of
modulus one its low update is
\begin{equation}
\label{eq:seed-echo-low}
 A_r\longmapsto W_Z(g_r)A_r,
 \qquad
 g_r=\frac{t}{r}(v'_r-v_r)=-2t w_r+o(1),
\end{equation}
where \(v_r\) is the actual \(Z\)-insertion before the forward cost
and \(v'_r\) is the actual insertion immediately after the selective
flip.  The special-center update can be taken to be exactly
\begin{equation}
\label{eq:seed-D-update}
 D_r\longmapsto D_r-t(v_r+v'_r).
\end{equation}
All ordinary errors are \(o(r)\) before and after the large costs.
\end{lemma}

\begin{proof}
Write \(u_{\mathrm{in},r}\) for the exact prefix before the echo,
and \(V_r\) for its selective flip. The next two prefixes are
\[
 u_{\mathrm{cost},r}=e^{-\ii tZq(v_r)/r}u_{\mathrm{in},r},
 \qquad u_{\mathrm{flip},r}=V_r u_{\mathrm{cost},r}.
\]
Their actual \(Z\)-insertions satisfy
\[
 \begin{aligned}
 u_{\mathrm{in},r}^*Zu_{\mathrm{in},r}\vac
 &=u_{\mathrm{cost},r}^*Zu_{\mathrm{cost},r}\vac=v_r,\\
 u_{\mathrm{flip},r}^*Zu_{\mathrm{flip},r}\vac&=v'_r.
 \end{aligned}
\]
The first equality holds because the forward cost commutes with the
root \(Z\). It splits the ordinary
part into spins \(y\) with centers \(-tyv_r\), and changes each
special center to \(d_{zy,r}-tyv_r\).  Centered tails remain bounded.
The standard-stage imaginary projection against \(v_r\) is unchanged
by adding \(-tyv_r\).  Thus the assumed special visibility separates
special and ordinary centers at this huge-cost stage.

Denote the exact ordinary and special components after that forward
cost by \(\Psi^{\mathrm{cost}}_{O,r}\) and
\(\psi^{\mathrm{cost}}_{S,r}\), respectively.
Lemma~\ref{lem:seed-selective-accuracy}
gives net identity to \(o(r)\) on the ordinary part, including its
\(Z\)-inserted version.  On the special part it gives \(-\ii X\)
to error \(o(r)\), again including the \(Z\)-inserted version.
Since \(XZX=-Z\), it follows that
\begin{align}
 v'_r-v_r
 &=u_{\mathrm{cost},r}^*(V_r^*ZV_r-Z)
      (\Psi^{\mathrm{cost}}_{O,r}(\vac)+\psi^{\mathrm{cost}}_{S,r})\notag\\
 &=-2u_{\mathrm{cost},r}^*Z\psi^{\mathrm{cost}}_{S,r}+o(r)
 =-2r w_r+o(r).
\label{eq:seed-flip-change}
\end{align}
The last equality uses transport, which the forward cost preserves.
The same calculation shows that transport changes sign at this flip.

On an ordinary input the reverse cost combines with the forward cost
as
\begin{equation}
\label{eq:seed-exact-echo-cancellation}
 e^{\ii tZq(v'_r)/r}e^{-\ii tZq(v_r)/r}
 =\xi_r W_Z\!\left(\frac{t}{r}(v'_r-v_r)\right),
 \qquad |\xi_r|=1.
\end{equation}
The Weyl phase \(\xi_r\) is independent of the root spin because the
two spin factors multiply to \(Z^2=\1\).  It may fail to converge.
Equation~\eqref{eq:seed-flip-change} makes the residual direction
bounded and proves~\eqref{eq:seed-echo-low}.  The ordinary \(o(r)\)
probe error is carried through the reverse cost by unitarity.  No
insertion approximation has been divided by \(r\): the cancellation
in~\eqref{eq:seed-exact-echo-cancellation} used the two actual
directions, and only then was~\eqref{eq:seed-flip-change} invoked.

A special packet has spin \(-y\) between the flips.  Its two cost
displacements therefore add as \(-tyv_r-tyv'_r\), giving
\eqref{eq:seed-D-update}.  Here \(y\) remains its label from the
standard stage before the echo; it is not relabeled when the spin
flips.  Its scalar phases are harmless.  The second
selective flip is applied at a standard ordinary geometry, since
\eqref{eq:seed-exact-echo-cancellation} has already removed the huge
ordinary centers.  It acts trivially to \(o(r)\) on ordinary tails,
restores special spin \(y\), and restores the plus sign in transport.
At this second-flip probe start the exact special norm is still \(r\),
and the ordinary low action already has the updated unitary
\(W_Z(g_r)A_r\).  Lemma~\ref{lem:seed-low-insertions} consequently
approximates the actual insertion to \(O(r)\) by this updated low
insertion, irrespective of the temporarily negative transport sign.

The special insertion remains transported because each cost commutes
with \(Z\), while each selective flip anticommutes with \(Z\) on the
special component to error \(o(r)\).  In particular, this proof needs
only forward estimates on that component.  Boundedness of the residual
direction implies uniform moment bounds for the new \(A_r,A_r^*\).
Lemma~\ref{lem:seed-low-insertions} supplies bounded approximants to
the next actual insertion.  Finally, \(v'_r-v_r=O(r)\), so the updated
\(D_r\) retains the asserted bounded-approximant property.
\end{proof}

This induction also proves its claimed uniformity on input families.
The exact decomposition~\eqref{eq:seed-exact-special-norm} is linear in
the input.  At initialization every input has the same four special
centers.  Unscaled costs and compensated mixers preserve that center
list, and~\eqref{eq:seed-D-update} updates it independently of the tail.
The probe parameters are determined on the vacuum, but
Lemma~\ref{lem:seed-selective-accuracy} acts uniformly on the whole
family of tails having these centers.  Every discarded vector has norm
\(o(r)\) and is propagated by exact unitaries.  Since the total gate
and probe lists are fixed and finite before \(r\downarrow0\), their
errors add to \(o(r)\).  This establishes the invariant with no
assertion about number moments of the discarded remainders.

\subsection{The external-mode limit and the energy pairing}

The echo has now produced the bounded direction \(-2t w_r+o(1)\).
It remains to show that this direction acts as an independent Gaussian
seed, and that the resulting identification preserves root-shift
energy.  Ordinary Gram convergence alone does not give the latter
conclusion, so we compute the root-shift pairing explicitly.

We now identify the limiting low circuit and at the same time establish
the limits of the projections used to select the probes.  Introduce at
each vertex a two-dimensional real external mode space with orthonormal
vectors \(\ell,w\), and let \(E_{\mathrm{ext}}=\Span_{\C}\{\ell,w\}\).
Its Gaussian position operators commute and are independent standard
Gaussians in the vacuum.  To specify the recursion by height, put
\[
 \cH_{\mathrm{seed},0}=\C^2\otimes\Gamma_s(E_{\mathrm{ext}}),
 \qquad
 \cH_{\mathrm{seed},h+1}
 =\C^2\otimes\Gamma_s\!\left(
    \cH_{\mathrm{seed},h}^{\circ}\oplus E_{\mathrm{ext}}\right).
\]
The vacuum is \(\ket+\otimes\Omega\), and empty descendants give the
same compatible height embeddings as in the foundations.  In the mode
space at height \(h+1\), the first summand contains the whole centered
child rooted space, including that child's seeds.  The second summand
contains the two new seeds at the current root.  Equivalently, the root
Hilbert space is the spin times the pointed Gaussian space of these two
seeds times the symmetric Fock space over centered children.

The formal low circuit starts with
\begin{equation}
\label{eq:seed-formal-initialization}
 A^{\mathrm{form}}=R_\star e^{2\ii bZq(\ell)}.
\end{equation}
A bulk mixer and an unscaled cost have their usual seeded rules.  An
echo with parameter \(t\) has the formal update
\(e^{-2\ii t Zq(w)}\). At each vertex, \(q(w)\) is the Gaussian
seed of the requested word: the same coordinate is reused at every
seed gate, while different vertices have independent copies.
The coordinate \(q(\ell)\) is independent of those seeds and enters
through initialization; its effect is removed by taking \(b\) small
in the final proof. Thus an echo with \(t=\sigma/2\) has the seed
law of the requested gate of angle \(\sigma\), not a newly sampled
seed at each pulse. The root-shift calculation below identifies the
energy, which is not determined by this marginal law alone.

At a standard stage write
\(v_P^{\mathrm{form}}=(A^{\mathrm{form}})^*P
A^{\mathrm{form}}\vac\), with the appropriate accumulated word in
\(A^{\mathrm{form}}\).  Treat these insertions and \(e=Z\vac\) as
endogenous child directions.  In particular their inner products with
the external directions \(\ell,w\) are zero.

\begin{lemma}[Finite Gram induction]
\label{lem:seed-gram-induction}
Suppose the required selective probe lists have been constructed through
a given standard stage.  The joint Gram matrix of \(e\), all actual
standard-stage insertion vectors used so far, and \(\ell_r,w_r\),
converges to that of the corresponding formal endogenous directions
and the two orthogonal external directions.  The limiting energy of the
actual insertion converges to the energy of the formal seeded insertion.
The same assertions hold at the standard ordinary geometry just before
the second flip in an echo, with the temporary transport sign retained.
\end{lemma}

\begin{proof}
Every \(A_r\) in the invariant is a finite sum of root spin matrices
times Weyl displacements.  Its directions are finite linear combinations
of earlier actual insertion directions, \(d_{0,r}\), and the bounded
echo directions \(g_r\).  Its adjoint and its insertion vectors have
the same finite-Weyl-sum form.  The number of summands is finite for the
fixed circuit; no bound uniform in circuit length is required.

For such sums all vacuum inner products reduce to
\begin{equation}
\label{eq:seed-vacuum-weyl-gram}
 \braket{W(f)\Omega}{W(g)\Omega}
 =\exp\!\left(-\tfrac12\norm{g-f}^2
                 +\ii\im\braket f g\right).
\end{equation}
Thus only finitely many direction inner products are needed at a step.
The initial directions satisfy \(d_{0,r}=2b\ell_r+o(1)\), and the
echo directions satisfy \(g_r=-2tw_r+o(1)\).
Norm-small errors in these bounded directions may be removed using
Lemma~\ref{lem:seed-packet-estimates}.  Every actual standard insertion
differs by \(O(r)\) from its bounded smooth low insertion, by
Lemma~\ref{lem:seed-low-insertions}.  Thus, when computing a new Gram
entry, an earlier actual insertion may first be replaced by its low
insertion at error \(O(r)\).  Lemma~\ref{lem:seed-high-directions}
therefore makes its pairing with \(\ell_r,w_r\) vanish.  Together with
\(\braket{\ell_r}{w_r}\to0\), this is exactly the orthogonal-sum rule
for external and endogenous directions.  Formula
\eqref{eq:seed-vacuum-weyl-gram} now proves the Gram assertion inductively.

Here is the additional calculation for the root-shift energy, so that
ordinary Gram convergence is not silently substituted for convergence
of a different operator.  A low \(Y\)- or \(Z\)-insertion is centered:
the circuit preserves the joint spin-flip parity and these matrices are
odd.  Write such an insertion as
\[
 h=\sum_{j=1}^J s_j\otimes W(f_j)\Omega,
 \qquad s_j\in\C^2.
\]
In the unseeded representation,
\(\cL h=\ket+\otimes a^*(h)\Omega\).  The one-particle coefficient
of \(W(f)\Omega\) is
\(\ii e^{-\norm f^2/2}f\).  Consequently
\begin{equation}
\label{eq:seed-energy-weyl-sum}
 \braket h{\cL h}
 =-\ii\sum_{j=1}^J
       \braket{s_j}{+}\,e^{-\norm{f_j}^2/2}\braket{f_j}h.
\end{equation}
Each pairing \(\braket{f_j}h\) is one of the same finite linear
combinations of insertion and external-direction Grams.  In the formal
seeded representation the same formula holds, with \(h\) inserted in
the endogenous summand; the external components of \(f_j\) then pair
to zero. More explicitly, the external two-mode Fock factor represents
the Gaussian variables at the root, while the child Fock factor is
built over endogenous directions. The root shift puts the root-seed
factor in its pointed vacuum and creates \(h\) only in that child
factor. The two-factor decomposition is identified with Fock space
over their orthogonal direct sum.  The high-direction argument gives
exactly this zero pairing in the actual limit.  Finally, the actual insertion differs in norm
from \(h\) by \(O(r)\).  Since \(\cL\) is an isometry,
\[
 \abs{\cE(v)-\cE(h)}
 \leq(\norm v+\norm h)\norm{v-h},
\]
which proves energy convergence.  A finite scalar bulk phase cancels
from every insertion.  The same argument applies just before an echo's
second flip, because its ordinary centers have already canceled and its
low update has already been made.
\end{proof}

The standard-stage Gram limits determine the projections needed for
ordinary probe starts.  For a probe immediately after a cost of size
\(t/r\), we also need the following transverse projection estimate.

\begin{lemma}[The huge-stage transverse projection]
\label{lem:seed-huge-transverse}
Immediately after a forward cost with fixed \(t\neq0\), the actual
\(Y\)-insertion has vanishing pairing with every bounded family in
\(\mathcal S\).  In particular, its pairing with every special scaled
center tends to zero.
\end{lemma}

\begin{proof}
For a bounded input \(\eta_r\), use the standard-stage uniform forward
approximation on both \(\vac\) and \(\eta_r\).  After the huge cost
their ordinary pieces have centers \(-tyv_r\) and root spins \(y\).
In the pairing defining
\(\braket{u_r^*Yu_r\vac}{\eta_r}\), the inserted \(Y\) reverses the
spin.  Nonzero spin pairings therefore require the two opposite
ordinary centers.  Their distance is \(2|t|\), so their contribution
vanishes by~\eqref{eq:seed-packet-separation}.  Every contribution with
a special component is \(O(r)\); the ordinary remainders are \(o(r)\).
This proves the first assertion.  Before the huge cost, the invariant
approximates every special center by a bounded smooth vector to
\(O(r)\).  In particular the intermediate directions \(v'_r\) in
\eqref{eq:seed-D-update} differ by \(O(r)\) from the corresponding
earlier standard insertions, by~\eqref{eq:seed-flip-change}.
Adding the forward displacement
\(-tyv_r\) preserves this property.  The first assertion and
Cauchy--Schwarz prove the second.
\end{proof}

\subsection{Visibility and the order of construction}

The preceding arguments are conditional only on the visibility
conditions.  We now choose the target angles and all finite probe lists
so that those conditions hold.  The order is important: first use the
formal low circuit to perturb the target angles, then construct the
probe lists successively from limits established at earlier stages,
and only afterwards choose \(r\).  Fix a finite seeded word, regard all
its angles as independent real variables \(\lambda\), and replace each
seed gate of angle \(\sigma\) formally by an echo parameter
\(t=\sigma/2\).  Include in a finite list all
standard-stage probe starts: compensated mixers and the second flips
of echoes.  Also include the standard stage preceding each forward huge
cost.  At a second-flip probe start the current echo's reverse cost has
already occurred: count that echo in the center accumulator and include
its seed displacement in the formal ordinary state.  Its packet label
\(y\) is still the label before the echo, although its current spin is
\(-y\).  Define, at each probe start,
\[
 D^{\mathrm{form}}
 =-2\sum_{j:\,\text{reverse cost already applied}}
       t_jv_{Z,j}^{\mathrm{form}},
\]
where the insertion is taken immediately before that echo.  At each
listed stage and for each \(z,y\), impose
\begin{equation}
\label{eq:seed-formal-visibility}
 a_{zy}^{\mathrm{form}}(\lambda)
 =-2\im\braket{v_Z^{\mathrm{form}}}
                   {-2bz e+yD^{\mathrm{form}}}\neq0.
\end{equation}

Each function in~\eqref{eq:seed-formal-visibility} is real analytic in
the finitely many angles.  One may verify this without an
unbounded-operator analyticity theorem: expand the finite word into
Weyl sums and use~\eqref{eq:seed-vacuum-weyl-gram} recursively.  These
expressions use finitely many sums, products, and exponentials of
previous finite Gram expressions.

None of these analytic functions is identically zero.  To see this,
set \emph{all} angles of the target word to zero, while keeping the
initialization fixed.  Every listed formal standard stage is then
\eqref{eq:seed-formal-initialization}, and
\(D^{\mathrm{form}}=0\).  A direct one-Gaussian calculation gives
\[
 \braket{v_Z^{\mathrm{form}}}e
 =\ii e^{-8b^2},\qquad
 a_{zy}^{\mathrm{form}}=4bz e^{-8b^2}\neq0.
\]
The temporary minus sign in special transport before a second flip
does not change this formal ordinary insertion or its center formula.
Thus the same test point works for every listed stage.

The complement of the finite union of these analytic zero sets, and
of the hyperplanes \(t_j=0\), is dense.  Perturb the target angles
arbitrarily little into that complement.  Energy changes arbitrarily
little because it is continuous by the same finite-Weyl-sum formulas.

We can now construct the actual pulses successively.  At a standard
probe start, Lemma~\ref{lem:seed-gram-induction} supplies the limits of
both special projection coordinates, and
\eqref{eq:seed-formal-visibility} makes their first coordinate nonzero.
Choose the finite ensemble-control list from these limiting pairs.
Lemma~\ref{lem:seed-selective-accuracy} implements it with the required
errors.  At a huge-stage first flip, the first coordinate is unchanged
by the forward displacement, while the second coordinate tends to zero
by Lemma~\ref{lem:seed-huge-transverse}.  Its limiting pairs are
therefore \((a_{zy}^{\mathrm{form}},0)\), again nonzero.  They too
admit the ensemble-control construction.  Continue the packet and Gram
inductions to the next stage.  The argument at the second flip uses the
already updated formal low circuit, with the sign of transport still
negative until that flip is undone.  In particular no pulse list is
chosen using a limit that presupposes its own implementation.

\begin{proof}[Proof of Theorem~\ref{thm:seed-simulation}]
Start with the requested finite seeded word and a tolerance.  For a
small fixed \(b>0\), prepend
\(R_\star e^{2\ii b Zq(\ell)}\) in the formal seeded calculus, with
\(\ell\) an additional independent Gaussian coordinate at every site.
As \(b\downarrow0\), its energy tends to that of the original seeded
word.  Indeed, finite bounded-direction Weyl words are strongly
continuous in this initial displacement, as are their insertion and
energy Grams.  At \(b=0\) the initial global mixer acts on the original
\(\ket+\) product state by a scalar, so it changes no expectation;
the extra Gaussian coordinate is then unused.

Choose \(b\) sufficiently small for this error, and perturb the target
angles sufficiently little that all the visibility conditions hold,
with another arbitrarily small energy error.  Fix the resulting finite
probe lists by the successive construction above.  With these lists
fixed, take \(r>0\) sufficiently small.  The packet induction and
Lemma~\ref{lem:seed-gram-induction} show that the ordinary word consisting
of the initialization, unscaled costs, compensated mixers, and seed
echoes has energy arbitrarily close to that of the perturbed formal
seeded word.  Summing the three errors proves the assertion.

Every gate used is a cost or mixer pulse, possibly with negative angle.
Toggled probes and all compensations are finite such words; consecutive
gates of the same kind can be combined, or zero-angle gates inserted
to obtain the stipulated alternating QAOA format.  After \(b\), the
perturbation, the finite probe lists, and \(r\) have been chosen, the
entire word and its angles are fixed.  Proposition~\ref{prop:tree-transfer}
therefore applies before any approximation parameter is changed.
This proves the claimed order of limits and parameter independence.
\end{proof}
